\scanpage{170}
\mspage{28}
\diagram{\begin{tikzcd}[ampersand replacement=\&]
c \arrow[r] \& a\times_{\unsure{x|a}}b\\
e \arrow[u,dashed] \arrow[r,dashed] \& d \arrow[u]
\end{tikzcd}}

, avec $x,a,b,x,c,d$ dans $\unsure{X|A}$, on peut compléter en un diagramme \uline{commutatif}, avec $e$ dans $X$. En d'autres termes, tout diagramme commutatif à traits pleins ci-dessous se complète en y ajoutant un ``objet initial'' $e$

% margin note written sideways
\reversemarginpar\margin{$\times$ Elle signifie que $X$ satisfait PS1, et \deleted{\unsure{que p.}} la condition \uline{PS$'2$} \deleted{\unsure{Pour}} Tout diagramme comm. \diagram{\begin{tikzcd}[ampersand replacement=\&,column sep=small] a\arrow[r,shift left=.5ex]\arrow[r,shift right=.5ex]\&b\arrow[r]\&x\end{tikzcd}} s'insère dans un diagr. comm. \diagram{\begin{tikzcd}[ampersand replacement=\&,column sep=tiny] y\arrow[r]\&a\arrow[r,shift left=.5ex]\arrow[r,shift right=.5ex]\&b\arrow[r]\&x\end{tikzcd}}}\normalmarginpar

\diagram{\begin{tikzcd}[ampersand replacement=\&]
\& c\arrow[r]\arrow[ddr]\arrow[rrd,"{\deleted{\unsure{\longrightarrow}}}" description,rightsquigarrow] \& a\arrow[dr] \& \\
e\arrow[ur,dotted]\arrow[dr,dotted] \& \& \& x\\
\& d\arrow[uur]\arrow[r] \& b\arrow[ur] \&
\end{tikzcd}}

\margin{il suffit d'exiger pour $e$ la commutativité des deux carrés $(e c b d)$ et $(e d a c)$}

Cette condition est donc impliquée par $X$ \inserted{loc.} cofiltrant, et elle est satisfaite aussi \inserted{(par ce qui précède)} \inserted{plus généralement} \unsure{ssi|si} les $X_{/x}$ sont cofiltrants.$^{\times}$ Mais la \unsure{réciproque|réciprocité} est-elle vraie -- cette condition de \unsure{totale} $0$-connexité des $X_{/a}$ implique-t-elle que $X$ \inserted{soit tot. $0$-connexe ?} \deleted{Mais} Je n'ai pas d'exemples où les $X_{/x}$ sont \inserted{ou seulement même des $X$ $0$-connexes, les $X_{/x}$ cofiltr., sans que $X$ \unsure{même} soit cofiltr.} sans que $X$ soit loc. cofiltrant ;

% slanted
\reversemarginpar\margin{si et seulement si ?}\normalmarginpar

% slanted
\reversemarginpar\margin{Mais il y a un \unsure{de|des} \uuline{contre-exple} évident, si $X$ est un \uline{groupoïde}, \unsure{or|on} \unsure{tous} les $X_{/x}$ sont $\simeq e$ donc cofiltrants, mais $X$ lui n'est \deleted{loc.} cofiltrant (\unsure{i.e.} \unsure{contractile}) que si $X$ est \unsure{équivalent|discret} \unsure{à $e$}. \unsure{$\pi_1\ne1$}. \unsure{Voir} ci. \unsure{Si ça est vrai|Si c'est vrai}, \unsure{on tient donc} \unsure{2.8 bis} \unsure{un} \unsure{tel} que les $X_{/x}$ sont cofiltrants, \unsure{Ctr ex. ?} mais $X$ \unsure{a un $\pi_1$|ne l'est} \unsure{non|pas} \unsure{trivial|cofiltrant}.}\normalmarginpar

\unsure{ou} d'exemples où les \deleted{\unsure{produit}} $X_{/x}$ sont $0$-connexes, \deleted{\unsure{tels}} que les $\unsure{X'_{/x}|X_{/x}}$ soient cofiltrants. Mais notons le

\uline{Corollaire} 2.12 : Les conditions équivalentes de \deleted{\unsure{2,}} a) a') b) d) [$\,$à l'exclusion de c, c', plus faibles apparemment$\,$] du th. 2.6.

% archive sheet number: 170
